\chapter{Criterio de auditoría mínima de claims}

\section{Auditoría mínima}
\begin{definition}
Se llama auditoría mínima de claims al conjunto disciplinado de verificaciones básicas mediante las cuales un claim registrado en Tomo III puede ser revisado en cuanto a origen, alcance, condiciones de validez y condiciones de fallo.
\end{definition}

\begin{observation}
La auditoría mínima no equivale todavía a validación exhaustiva ni a cierre probatorio fuerte. Su función es fijar un primer umbral de revisión suficiente para evitar opacidad, sobrerreclamo o pérdida de trazabilidad.
\end{observation}

\section{Condición mínima de auditabilidad}
\begin{definition}
Un claim es mínimamente auditable cuando pueden declararse sin contradicción inmediata:
\begin{enumerate}
  \item su demo o fuente de origen;
  \item su dependencia canónica;
  \item su alcance declarado;
  \item sus condiciones mínimas de validez;
  \item sus condiciones mínimas de fallo;
  \item la evidencia o trazas que sostienen su formulación disciplinada.
\end{enumerate}
\end{definition}

\begin{observation}
No se exige todavía una teoría completa de evidencia final. Basta con una revisión suficiente para determinar si el claim sigue siendo legible, disciplinado y metodológicamente defendible.
\end{observation}

\section{Auditoría y tabla de claims}
\begin{definition}
La tabla de claims debe servir como soporte mínimo de auditoría, permitiendo revisar cada claim mediante una ficha común y comparable.
\end{definition}

\begin{observation}
La tabla no reemplaza la evidencia primaria, pero sí permite localizar con rapidez qué afirmación se audita, desde qué demo proviene y bajo qué alcance pretende sostenerse.
\end{observation}

\section{Auditoría y sobriedad}
\begin{definition}
La auditoría mínima exige distinguir entre:
\begin{enumerate}
  \item lo que la evidencia sostiene directamente;
  \item lo que la formulación interpreta de manera disciplinada;
  \item lo que solo puede mantenerse como proyección o hipótesis de trabajo.
\end{enumerate}
\end{definition}

\begin{observation}
Esta distinción protege al Tomo III frente a claims inflados que mezclen soporte efectivo con extrapolación no declarada.
\end{observation}

\section{Condición mínima de observación crítica}
\begin{definition}
Un claim debe marcarse como críticamente observado cuando, aun siendo registrable, presenta fragilidad visible en origen, alcance, trazabilidad o relación con su evidencia de apoyo.
\end{definition}

\begin{observation}
La observación crítica no obliga a expulsar el claim de inmediato. Su función es impedir que una afirmación frágil circule como si ya estuviera estabilizada.
\end{observation}

\section{Condición mínima de no auditabilidad}
\begin{definition}
Se considera que un claim no es todavía mínimamente auditable cuando no puede reconstruir de forma legible su origen, su alcance, sus condiciones de validez o la relación entre evidencia y formulación.
\end{definition}

\begin{observation}
Esta condición protege al corpus frente a afirmaciones que todavía no han alcanzado el grado mínimo de revisión metodológica exigible en Tomo III.
\end{observation}

\section{Cláusula de alcance}
\begin{observation}
El criterio de auditoría mínima de claims fija una primera capa de revisión metodológica dentro de Tomo III. No constituye todavía una auditoría final completa ni una validación exhaustiva del conjunto del corpus.
\end{observation}

